\section{2022 Applied \& Computational Mathematics}

\begin{exercise}
\label{applied:2022:1}
Consider $\left\{ p _ { i } ( x ) \right\} _ { i = 0 } ^ { \infty }$ , a family of orthogonal polynomials associated with the inner product

\[
\langle f, g \rangle = \int_ {- 1} ^ {1} f (x) g (x) w (x) d x, \quad w (x) > 0 \quad \mathrm {f o r} x \in (- 1, 1),
\]

where $p _ { i } ( x )$ is a polynomial of degree ??. Let $x _ { 0 } , x _ { 1 } , \ldots , x _ { n }$ be the roots of $p _ { n + 1 } ( x )$ . Construct an orthonormal basis in the subspace of the polynomials of degree no more than $n$ such that, for any polynomial in this subspace, the coefficients of its expansion into the basis are equal to the scaled values of this polynomial at the nodes $x _ { 0 } , x _ { 1 } , \ldots , x _ { n }$ .
\end{exercise}

\begin{solution}
Let $V = \mathbb{P}_n$ be the space of polynomials of degree at most $n$. The roots $x_0, x_1, \dots, x_n$ of the orthogonal polynomial $p_{n+1}(x)$ are the nodes for the $(n+1)$-point Gaussian quadrature rule associated with the weight function $w(x)$. This quadrature rule is exact for polynomials of degree at most $2n+1$.

Let $l_i(x)$ for $i=0, \dots, n$ be the Lagrange interpolating polynomials of degree $n$ for the nodes $x_0, \dots, x_n$:
\[ l_i(x) = \prod_{j \neq i}^n \frac{x-x_j}{x_i-x_j}. \]
Since $l_i(x)l_j(x)$ has degree $2n \leq 2n+1$, the Gaussian quadrature is exact:
\[ \langle l_i, l_j \rangle = \int_{-1}^{1} l_i(x)l_j(x)w(x) \, dx = \sum_{k=0}^{n} w_k l_i(x_k) l_j(x_k), \]
where $w_k$ are the positive Gaussian weights. Using the property $l_i(x_k) = \delta_{ik}$, we have:
\[ \langle l_i, l_j \rangle = \sum_{k=0}^{n} w_k \delta_{ik} \delta_{jk} = w_i \delta_{ij}. \]
Thus, the Lagrange polynomials are orthogonal. We normalize them to obtain an orthonormal basis $\{R_i\}_{i=0}^n$:
\[ R_i(x) = \frac{1}{\sqrt{w_i}} l_i(x). \]
For any polynomial $f(x) \in \mathbb{P}_n$, its expansion in this basis is $f(x) = \sum_{i=0}^n f_i R_i(x)$. The coefficients $f_i$ are given by the projections:
\[ f_i = \langle f, R_i \rangle = \int_{-1}^{1} f(x) \frac{l_i(x)}{\sqrt{w_i}} w(x) \, dx. \]
Since $f(x)l_i(x)$ has degree $2n \leq 2n+1$, Gaussian quadrature is again exact:
\[ f_i = \sum_{k=0}^n w_k f(x_k) \frac{l_i(x_k)}{\sqrt{w_i}} = \sum_{k=0}^n w_k f(x_k) \frac{\delta_{ik}}{\sqrt{w_i}} = \sqrt{w_i} f(x_i). \]
Thus, the coefficients $f_i$ are precisely the values of the polynomial at the nodes $x_i$ scaled by $\sqrt{w_i}$.
\end{solution}

\begin{exercise}
\label{applied:2022:2}
Consider a 2D fixed point iteration of the form

\[
x _ {k + 1} = f \left(x _ {k}, y _ {k}\right), y _ {k + 1} = g \left(x _ {k}, y _ {k}\right). \tag {1}
\]

Assume that the vector-valued function $\vec { H } ( x , y ) = ( f ( x , y ) , g ( x , y ) ) ^ { T }$ is continuously-differentiable, and the infinity norm of the Jacobian matrix is less than 1 at a unique fixed point $( x _ { \infty } , y _ { \infty } )$ .

Now consider a new iteration:

\[
x _ {k + 1} = f \left(x _ {k}, y _ {k}\right), \quad y _ {k + 1} = g \left(x _ {k + 1}, y _ {k}\right). \tag {2}
\]

Prove that iteration (2) is convergent, to the same fixed point as iteration (1), for the initial conditions sufficiently close to the fixed point.
\end{exercise}

\begin{solution}
First, we verify that the new iteration has the same fixed point as the original. For (1), the fixed point satisfy $x_\infty = f(x_\infty, y_\infty)$ and $y_\infty = g(x_\infty, y_\infty)$. Substituting into (2):
\[ f(x_\infty, y_\infty) = x_\infty, \quad g(x_\infty, y_\infty) = y_\infty. \]
Since $x_{k+1}$ in $g(x_{k+1}, y_k)$ becomes $x_\infty$ at the fixed point, the equations are identical, and $(x_\infty, y_\infty)$ is indeed a fixed point of iteration (2).

Next, we analyze the local convergence. Let $\mathbf{z}_k = (x_k, y_k)^T$. Iteration (2) can be seen as $\mathbf{z}_{k+1} = \Phi(\mathbf{z}_k)$ where $\Phi(x, y) = (f(x, y), g(f(x, y), y))^T$. The Jacobian of $\Phi$ at the fixed point is:
\[ \mathbf{J}_2 = \begin{pmatrix} f_x & f_y \\ g_x f_x & g_x f_y + g_y \end{pmatrix}. \]
The infinity norm of a matrix is the maximum absolute row sum. 
For the first row, the sum is $|f_x| + |f_y|$, which is the first row sum of the original Jacobian $\mathbf{J}_1$. By assumption:
\[ |f_x| + |f_y| \leq \|\mathbf{J}_1\|_\infty < 1. \]
For the second row, we have:
\begin{align*}
|g_x f_x| + |g_x f_y + g_y| &\leq |g_x| (|f_x| + |f_y|) + |g_y| \\
&\leq |g_x| \|\mathbf{J}_1\|_\infty + |g_y| \\
&< |g_x| + |g_y| \leq \|\mathbf{J}_1\|_\infty < 1.
\end{align*}
Thus, $\|\mathbf{J}_2\|_\infty < 1$, which implies the spectral radius $\rho(\mathbf{J}_2) \leq \|\mathbf{J}_2\|_\infty < 1$. By the Ostrowski Theorem, the iteration (2) converges locally to $(x_\infty, y_\infty)$.
\end{solution}

\begin{exercise}
\label{applied:2022:3}
Let $A \in \mathbf { R } ^ { \mathbf { m } \times \mathbf { m } }$ be a matrix with entries $a _ { i j }$ which satisfy

\[
a _ {i i} \geq \sum_ {j \neq i} | a _ {i j} | + 2, \quad a _ {i i} \leq 7.
\]

(a) Prove that $A ^ { - 1 }$ exists.   
(b) Prove that $\| A \| _ { \infty }$ is the max row sum (of absolute values) of $A$ .   
(c) Find both a lower and upper bound for $\| A \| _ { \infty }$   
(d) Now assume $A = A ^ { T }$ . Find bounds for $\| A \| _ { 2 }$ and $\| A ^ { - 1 } \| _ { 2 }$ .
\end{exercise}

\begin{solution}
\textbf{(a)} Suppose $A$ is singular. Then there exists a non-zero vector $v$ such that $Av = 0$. Let $|v_k| = \max_j |v_j| > 0$. From the $k$-th row of $Av=0$:
\[ a_{kk} v_k = -\sum_{j \neq k} a_{kj} v_j \implies |a_{kk}| |v_k| \leq \sum_{j \neq k} |a_{kj}| |v_j| \leq \left(\sum_{j \neq k} |a_{kj}|\right) |v_k|. \]
Dividing by $|v_k|$, we get $|a_{kk}| \leq \sum_{j \neq k} |a_{kj}|$. This contradicts the given condition $a_{ii} \geq \sum_{j \neq i} |a_{ij}| + 2$ (which implies $a_{ii} > \sum_{j \neq i} |a_{ij}|$). Thus $A$ must be non-singular, and $A^{-1}$ exists.

\textbf{(b)} The matrix infinity norm is $\|A\|_\infty = \sup_{x \neq 0} \frac{\|Ax\|_\infty}{\|x\|_\infty}$. For any $x$ with $\|x\|_\infty = 1$, we have:
\[ |(Ax)_i| = \left|\sum_j a_{ij} x_j\right| \leq \sum_j |a_{ij}| |x_j| \leq \sum_j |a_{ij}|. \]
Thus $\|Ax\|_\infty \leq \max_i \sum_j |a_{ij}|$.
Equality is achieved by choosing $k$ such that $\sum_j |a_{kj}|$ is maximal and setting $x_j = \text{sgn}(a_{kj})$. Then $\|x\|_\infty = 1$ and $|(Ax)_k| = \sum_j |a_{kj}|$. Thus $\|A\|_\infty = \max_i \sum_j |a_{ij}|$.

\textbf{(c)} Let $R_i = \sum_{j \neq i} |a_{ij}|$. The $i$-th row absolute sum is $a_{ii} + R_i$ (since $a_{ii} \geq 2$).
From $R_i \leq a_{ii} - 2$, we have:
\[ \|A\|_\infty = \max_i (a_{ii} + R_i) \leq \max_i (a_{ii} + a_{ii} - 2) = 2a_{ii} - 2 \leq 2(7) - 2 = 12. \]
For the lower bound, since $R_i \geq 0$, we have $a_{ii} + R_i \geq a_{ii} \geq 2$. Thus $2 \leq \|A\|_\infty \leq 12$.

\textbf{(d)} If $A = A^T$, then $\|A\|_2 = \rho(A)$. By Gershgorin's Disk Theorem, each eigenvalue $\lambda$ lies in some disk $D_i = [a_{ii} - R_i, a_{ii} + R_i]$.
As shown in (c), $a_{ii} - R_i \geq 2$ and $a_{ii} + R_i \leq 12$. Thus $2 \leq \lambda \leq 12$ for all eigenvalues.
Consequently, $2 \leq \|A\|_2 \leq 12$.
For $A^{-1}$, the eigenvalues are $1/\lambda$, so $1/12 \leq \|A^{-1}\|_2 \leq 1/2$.
\end{solution}

\begin{exercise}
\label{applied:2022:4}
Consider a system of ODE initial value problems of the form:

\[
\frac {d}{d t} u = f (u), \quad u (0) = u _ {0}.
\]

Assume that $f ( u )$ has the property that the forward Euler (FE) method:

\[
U ^ {n + 1} = U ^ {n} + k f (U ^ {n}),
\]

satisfies

\[
\| U ^ {n + 1} \| \leq \| U ^ {n} \|
\]

for some norm $\| \cdot \|$ and for all time-steps $k$ , $0 < k \leq k _ { F E }$ . Now consider the 2-stage Runge-Kutta method:

\[
\begin{array}{l} U ^ {(1)} = U ^ {n} + k \beta_ {1 0} f (U ^ {n}), \\ U ^ {n + 1} = \left\{\alpha_ {2 0} U ^ {n} + k \beta_ {2 0} f (U ^ {n}) \right\} + \left\{\alpha_ {2 1} U ^ {(1)} + k \beta_ {2 1} f (U ^ {(1)}) \right\} \\ \end{array}
\]

where

\[
\beta_ {1 0} \geq 0, \quad \beta_ {2 0} \geq 0, \quad \beta_ {2 1} \geq 0, \quad \alpha_ {2 0} \geq 0, \quad \alpha_ {2 1} \geq 0, \quad \alpha_ {2 0} + \alpha_ {2 1} = 1.
\]

(a) Prove that the above 2-stage Runge-Kutta method also satisfies the inequality:

\[
\| U ^ {n + 1} \| \leq \| U ^ {n} \|
\]

under some appropriate time-step restriction: $0 \leq k \leq k ^ { * }$ , where you need to explicitly determine $k ^ { * }$ in terms of $k _ { F E }$ .

(b) Explicitly determine the coefficients:

\[
\beta_ {1 0}, \quad \beta_ {2 0}, \quad \beta_ {2 1}, \quad \alpha_ {2 0}, \quad \alpha_ {2 1},
\]

so that

(i) The method is second-order accurate; and   
(ii) The maximum allowed time-step, $k ^ { * }$ , is as large as possible.
\end{exercise}

\begin{solution}
\textbf{(a)} The stage $U^{(1)}$ is a Forward Euler step with step size $k\beta_{10}$. It is stable ($\|U^{(1)}\| \leq \|U^n\|$) if $k\beta_{10} \leq k_{\mathrm{FE}}$, i.e., $k \leq k_{\mathrm{FE}}/\beta_{10}$.
The final step is:
\[ U^{n+1} = \alpha_{20} \left(U^n + \frac{k\beta_{20}}{\alpha_{20}} f(U^n)\right) + \alpha_{21} \left(U^{(1)} + \frac{k\beta_{21}}{\alpha_{21}} f(U^{(1)})\right). \]
This is a convex combination of two Forward Euler steps from $U^n$ and $U^{(1)}$ respectively. They are stable if:
\[ \frac{k\beta_{20}}{\alpha_{20}} \leq k_{\mathrm{FE}} \quad \text{and} \quad \frac{k\beta_{21}}{\alpha_{21}} \leq k_{\mathrm{FE}}. \]
Under these conditions, $\|U^{n+1}\| \leq \alpha_{20} \|U^n\| + \alpha_{21} \|U^{(1)}\| \leq (\alpha_{20} + \alpha_{21}) \|U^n\| = \|U^n\|$.
Thus $k \leq k^* = k_{\mathrm{FE}} \cdot \min \left( \frac{1}{\beta_{10}}, \frac{\alpha_{20}}{\beta_{20}}, \frac{\alpha_{21}}{\beta_{21}} \right)$.

\textbf{(b)} To evaluate accuracy, we apply the method to $u' = \lambda u$ (let $z = k\lambda$):
$U^{(1)} = (1 + z\beta_{10}) U^n$.
$U^{n+1} = \alpha_{20}(1 + z\frac{\beta_{20}}{\alpha_{20}})U^n + \alpha_{21}(1 + z\frac{\beta_{21}}{\alpha_{21}})U^{(1)}$.
Substituting $U^{(1)}$:
\begin{align*}
U^{n+1} &= \left[ (\alpha_{20} + z\beta_{20}) + (\alpha_{21} + z\beta_{21})(1 + z\beta_{10}) \right] U^n \\
&= \left[ (\alpha_{20}+\alpha_{21}) + z(\beta_{20} + \beta_{21} + \alpha_{21}\beta_{10}) + z^2(\beta_{10}\beta_{21}) \right] U^n.
\end{align*}
Matching with the Taylor series $e^z = 1 + z + z^2/2 + O(z^3)$:
1) $\alpha_{20} + \alpha_{21} = 1$ (given).
2) $\beta_{20} + \beta_{21} + \alpha_{21}\beta_{10} = 1$.
3) $\beta_{10}\beta_{21} = 1/2$.
To maximize $k^*$, we want to maximize $\min \left( \frac{1}{\beta_{10}}, \frac{\alpha_{20}}{\beta_{20}}, \frac{\alpha_{21}}{\beta_{21}} \right)$.
The optimal coefficients (SSP-RK2) are:
$\beta_{10} = 1, \beta_{21} = 1/2, \alpha_{21} = 1/2 \implies \beta_{20} = 0, \alpha_{20} = 1/2$.
Checking $k^*$: $\min(1/1, \infty, 1) = 1$. So $k^* = k_{\mathrm{FE}}$.
\end{solution}

\begin{exercise}
\label{applied:2022:5}
Construct a third-order accurate Lax-Wendroff-type method for $u _ { t } + a u _ { x } = 0$ $[ a > 0$ is a constant) in the following way:

(a) • Expand $u ( t + k , x )$ in a Taylor series and keep the first four terms. Replace all time derivatives by spatial derivatives using the equation.   
• Construct a cubic polynomial passing through the points $U _ { j - 2 } ^ { n } , U _ { j - 1 } ^ { n } , U _ { j } ^ { n } , U _ { j + 1 } ^ { n } .$   
• Approximate the spatial derivatives in the Taylor series by the exact derivatives of the above constructed cubic polynomial.

(b) Verify that the truncation error is $O ( k ^ { 3 } )$ if $h = O ( k )$ .
\end{exercise}

\begin{solution}
\textbf{(a)} Taylor expansion in time:
\[ u(t+k, x) = u + k u_t + \frac{k^2}{2} u_{tt} + \frac{k^3}{6} u_{ttt} + O(k^4). \]
From $u_t = -au_x$, we have $u_{tt} = a^2 u_{xx}$ and $u_{ttt} = -a^3 u_{xxx}$. Thus:
\[ u(t+k, x) = u - ak u_x + \frac{a^2k^2}{2} u_{xx} - \frac{a^3k^3}{6} u_{xxx} + O(k^4). \]
Let $p_3(x)$ be the cubic interpolating $\{U_{j-2}, U_{j-1}, U_j, U_{j+1}\}$ at $x_i = ih$. Its derivatives at $x_j$ are:
\begin{align*}
u_x &\approx p_3'(x_j) = \frac{U_{j-2} - 6U_{j-1} + 3U_j + 2U_{j+1}}{6h}, \\
u_{xx} &\approx p_3''(x_j) = \frac{U_{j-1} - 2U_j + U_{j+1}}{h^2}, \\
u_{xxx} &\approx p_3'''(x_j) = \frac{-U_{j-2} + 3U_{j-1} - 3U_j + U_{j+1}}{h^3}.
\end{align*}
The scheme is:
\[ U_j^{n+1} = U_j^n - \frac{ak}{6h}(U_{j-2} - 6U_{j-1} + 3U_j + 2U_{j+1}) + \frac{a^2k^2}{2h^2}(U_{j-1} - 2U_j + U_{j+1}) - \frac{a^3k^3}{6h^3}(-U_{j-2} + 3U_{j-1} - 3U_j + U_{j+1}). \]

\textbf{(b)} The local truncation error is $\tau = \frac{u(t+k) - u(t)}{k} - (\dots)$.
The spatial derivative approximations satisfy $p_3^{(m)} = u^{(m)} + O(h^{4-m})$.
Thus $p_3' = u_x + O(h^3)$, $p_3'' = u_{xx} + O(h^2)$, $p_3''' = u_{xxx} + O(h)$.
Substituting into the scheme:
\[ u(t+k) = u - ak(u_x + O(h^3)) + \frac{a^2k^2}{2}(u_{xx} + O(h^2)) - \frac{a^3k^3}{6}(u_{xxx} + O(h)) + O(k^4). \]
The error is:
\[ E = -ak O(h^3) + k^2 O(h^2) - k^3 O(h) + O(k^4). \]
If $h = O(k)$, then $E = O(k^4)$. The truncation error per step is $E/k = O(k^3)$.
\end{solution}

\begin{exercise}
\label{applied:2022:6}
Suppose you have $\$ 60 K$ to invest and there are 3 investment options available. You must invest in multiples of $\$ 10K$ . If $d _ { i }$ dollars are invested in investment ?? then you receive a net value (as the profit) of $r _ { i } ( d _ { i } )$ dollars. For $d _ { i } > 0$ we have

\[
\begin{array}{l} r _ {1} (d _ {1}) = (7 d _ {1} + 2) \times 1 0, \\ r _ {2} (d _ {2}) = (3 d _ {2} + 7) \times 1 0, \\ r _ {3} (d _ {3}) = (4 d _ {3} + 5) \times 1 0, \\ \end{array}
\]

and $d _ { 1 } ( 0 ) = d _ { 2 } ( 0 ) = d _ { 3 } ( 0 )$ . All are measured in $\$ 10K$ dollars. The objective is to maximize the net value of your

investments. This can be formulated as a linear programming problem:

\[
\max _ {d _ {1}, d _ {2}, d _ {3}} r _ {1} (d _ {1}) + r _ {2} (d _ {2}) + r _ {3} (d _ {3}),
\]

\[
\text {s u c h t h a t} \quad d _ {1} + d _ {2} + d _ {3} \leq 6,
\]

\[
d _ {i} \geq 0 \quad i = 1, 2, 3 \quad \text {a r e i n t e g e r s}.
\]
\end{exercise}

\begin{solution}
We solve this via dynamic programming. Let $f_i(s)$ be the max return from options $i$ to 3 with total investment $s$.
The returns $r_i(d)$ for $d \in \{0, \dots, 6\}$ are:
$r_1: [0, 90, 160, 230, 300, 370, 440]$
$r_2: [0, 100, 130, 160, 190, 220, 250]$
$r_3: [0, 90, 130, 170, 210, 250, 290]$

Stage 3: $f_3(s) = r_3(s)$.
Stage 2: $f_2(s) = \max_{0 \leq d \leq s} \{r_2(d) + f_3(s-d)\}$.
$f_2(1) = \max(r_2(1)+f_3(0), r_2(0)+f_3(1)) = \max(100+0, 0+90) = 100$ ($d_2=1$).
$f_2(2) = \max(r_2(2)+0, r_2(1)+90, 0+130) = \max(130, 190, 130) = 190$ ($d_2=1, d_3=1$).
$f_2(3) = \max(r_2(3), r_2(2)+90, r_2(1)+130, 170) = \max(160, 220, 230, 170) = 230$ ($d_2=1, d_3=2$).
... Following this, $f_2(6)=350$ (with $d_2=1, d_3=5$).
Stage 1: $f_1(6) = \max_{0 \leq d_1 \leq 6} \{r_1(d_1) + f_2(6-d_1)\}$.
$d_1=0: 0 + 35 = 35$
$d_1=1: 9 + 31 = 40$
$d_1=2: 16 + 27 = 43$
$d_1=3: 23 + 23 = 46$
$d_1=4: 30 + 19 = 49$
$d_1=5: 37 + 10 = 47$
$d_1=6: 44 + 0 = 44$
Maximum return is 490 (units of $\$1K$), achieved at $d_1=4, d_2=1, d_3=1$.
\end{solution}
